\documentclass{../../../unipd-notes}
\unipdsetup{course = {Fondamenti di Ingegneria Informatica},short-course = {Fondamenti di Ingegneria},document-type = {Appunti delle lezioni}}
\begin{document}
\tableofcontents
\chapter{Rappresentazione dell'informazione}
\section{Sistemi di numerazione}
Un numero binario $b_{n-1}\dots b_0$ rappresenta il valore $\sum_{i=0}^{n-1}b_i2^i$.
\subsection{Conversione in base dieci}
La parola $101101_2$ corrisponde a $45_{10}$.
\subsubsection{Verifica del risultato}
La riconversione in base due permette di controllare eventuali errori di calcolo.
\courseappendices
\chapter{Tabella dei codici ASCII}
L'appendice raccoglie i codici più utilizzati durante le esercitazioni.
\end{document}
